\documentclass[../main.tex]{subfiles}
\begin{document}

\FloatBarrier
\section{Overview}
\label{sec:overview}

\begin{figure}[H]
\centering
\begin{subfigure}{\linewidth}
\centering
\includegraphics[width=0.7\linewidth]{overview_baseline.png}
\caption{Single-encoder baselines.}
\label{fig:overview_baseline}
\end{subfigure}\\[2.5ex]
\begin{subfigure}{\linewidth}
\centering
\includegraphics[width=\linewidth]{overview_proposed.png}
\caption{Proposed feature-group hybrid model (detailed in Figure~\ref{fig:hybrid}).}
\label{fig:overview_hybrid}
\end{subfigure}
\caption{The baseline and the proposed hybrid model.}
\label{fig:overview_models}
\end{figure}

Two public datasets form the training data, while a self-collected set recorded on real webcams serves as an in-the-wild test set (Section~\ref{sec:data}). Every clip carries one of four ordinal engagement levels, from highly disengaged to highly engaged. Each clip is represented by per-frame OpenFace descriptors and a clip-level I3D embedding extracted from the raw video, then resampled to a common length and standardised with training-set statistics so that no target information leaks into preprocessing.

Two model paradigms are compared, contrasted in Figure~\ref{fig:overview_models}. The single-encoder baselines follow the standard approach, in which one encoder, an MLP, TCN, LSTM, or Transformer, consumes the full 709-dimensional OpenFace sequence, or an MLP consumes the pooled I3D vector, before a classifier produces the prediction (Section~\ref{sec:baselines}).

The proposed hybrid model replaces that single encoder with a split-and-encode design: the OpenFace descriptor is split into five meaning-based feature groups, each encoded independently by its own temporal encoder, the I3D vector is encoded by an MLP, and the resulting embeddings are concatenated and classified by a shared MLP while an auxiliary head supervises every stream (Section~\ref{sec:hybrid}). The loss functions, training regime, and 3$\times$3 cross-dataset protocol under which every model is trained and evaluated are established with the experiments that use them in Chapter~\ref{chap:results}.

\FloatBarrier
\section{Data and Preprocessing}
\label{sec:data}

\subsection{Public datasets}

Two public datasets supply the labelled training data. DAiSEE~\cite{daisee} is a widely used engagement dataset of short webcam clips, recorded in the wild, each showing a single learner under uncontrolled conditions and annotated on a four-level ordinal engagement scale, which makes it realistic but noisy; it ships official Train, Validation, and Test partitions used here directly for training, early-stopping validation, and testing. CMOSE~\cite{cmose} is a more recent and larger dataset, likewise collected in the wild but annotated by experts and reporting strong inter-rater agreement, and it ships precomputed OpenFace and I3D features; it is divided into a labelled training split, an \texttt{unlabel} split repurposed as the validation set for early stopping and checkpoint selection, and a test split. Both label every clip on the same four ordinal levels used throughout the thesis, E0 (Highly\nobreakdash-Disengage), E1 (Disengage), E2 (Engage), and E3 (Highly\nobreakdash-Engage), so that CMOSE, with its greater scale and label quality, serves as the primary source and DAiSEE as an independent second source.
\begin{table}[ht]
\centering
\caption{Dataset statistics and class balance (\%).}
\label{tab:T3_1_dataset_stats}
\fittable{%
\begin{tabular}{lrrrrrrrr}
\toprule
Dataset & Total & Train & Val & Test & E0\% & E1\% & E2\% & E3\% \\
\midrule
CMOSE & 12197 & 8783 & 2193 & 1221 & 2.800 & 18.100 & 69.400 & 9.600 \\
DAiSEE & 8571 & 5358 & 1429 & 1784 & 0.700 & 5.100 & 50.300 & 43.800 \\
\bottomrule
\end{tabular}%
}
\end{table}


Table~\ref{tab:T3_1_dataset_stats} shows that both public datasets are heavily imbalanced toward the upper-middle engagement levels, and imbalanced in different directions: CMOSE concentrates on the single Engage level, whereas DAiSEE divides more evenly between Engage and Highly\nobreakdash-Engage. The disengaged levels are rare in both, yet they are precisely the responses a monitoring system must catch. The same skew holds within each individual split, so it reflects the nature of the task rather than an artefact of how either dataset was partitioned. The self-collected private test set in the final row, though gathered independently, carries the same engaged-heavy imbalance, so the skew is a property of the engagement task and not of any single~collection.

\subsection{Private dataset}
\label{sec:private-collection}

The private set exists to measure deployment behaviour rather than in-dataset behaviour. No model is trained or selected on it, so it is a test-only set, and it is sampled from the webcam recordings a deployed model would actually face, so it probes a condition the public datasets leave untested. The reason this matters is geometric: the OpenFace gaze, head-pose, and landmark features are read from the camera's view of the face, so the same behaviour recorded from a different camera angle produces different feature values. Figure~\ref{fig:private-angles} shows private clips that differ along several axes: the camera's height and distance relative to the student's eye level, its position relative to the screen, its orientation, whether vertical or horizontal, and the background behind the learner, whether the real room, blurred, or a green screen. CMOSE and DAiSEE were recorded under far more uniform setups and so cannot exercise this variation, which makes a small, locally collected, hand-labelled set drawn from realistic online-class recordings the most direct way to check that a model stays fit once the camera geometry changes.

\begin{figure}[htbp]
\centering
\begin{subfigure}{0.48\linewidth}
\centering
\includegraphics[height=3.2cm]{private1.png}
\caption{}
\label{fig:private-a}
\end{subfigure}\hfill
\begin{subfigure}{0.48\linewidth}
\centering
\includegraphics[height=3.2cm]{private2.png}
\caption{}
\label{fig:private-b}
\end{subfigure}\\[1.5ex]
\begin{subfigure}{0.48\linewidth}
\centering
\includegraphics[height=3.2cm]{private3.png}
\caption{}
\label{fig:private-c}
\end{subfigure}\hfill
\begin{subfigure}{0.48\linewidth}
\centering
\includegraphics[height=3.2cm]{private4.png}
\caption{}
\label{fig:private-d}
\end{subfigure}
\caption{Representative frames from the private set, illustrating the range of camera angles found in unconstrained online-class recordings.}
\label{fig:private-angles}
\end{figure}

The set was built with the same feature pipeline as the public data. Webcam recordings of online-class sessions were sampled at several points and cut into fixed-length ten-second clips, each containing a single learner, and every clip was run through the OpenFace pipeline used for the public datasets, with clips that retained too few tracked frames discarded as unreliable. The accepted clips were then labelled through a purpose-built local interface, one engagement level per clip and by a single annotator, a reliance revisited as a limitation in Section~\ref{sec:summary}. Finally, the same OpenFace (709-d per frame) and I3D (1024-d per clip) extractors used for DAiSEE were applied, so the private features share the format of the precomputed CMOSE features and are distributionally comparable to the training data, and the 366 clips with complete, aligned OpenFace and I3D features form the final set. Its label prior differs again from both public datasets and, together with its independent recording conditions, makes it an in-the-wild measure of deployment behaviour that anchors Section~\ref{sec:private-result}.

\subsection{Preprocessing}

Each clip is represented by two feature streams. OpenFace~2.0~\cite{openface2} produces a 709-dimensional descriptor for every frame, and I3D~\cite{i3d} produces a single 1024-dimensional embedding pooled over the whole clip. CMOSE ships both precomputed; for DAiSEE and the private set they are extracted from the raw video with the same toolkits, so all three sources share one feature format. When an OpenFace frame reports more than one detected face the highest-confidence row is kept, and any non-numeric entry is set to zero. Because clips differ in length, each OpenFace sequence is resampled along time to a fixed number of frames by per-feature linear interpolation, 300 frames by default, reduced to 150 for the largest training runs to fit memory and to 75 in the multimodal variant, while the I3D stream, being a single vector per clip, needs no temporal resampling and is encoded by an MLP (Section~\ref{sec:hybrid}).

Every feature is then standardised by a per-feature $z$-score. The mean and standard deviation are fit on the training split alone, and the validation, test, and private sets are transformed with those training statistics, so no target-set information leaks into preprocessing. The OpenFace and I3D streams are standardised with separate~statistics.

\FloatBarrier
\section{Baseline Models}
\label{sec:baselines}

Five single-encoder baselines establish the comparison points, and their architectures are shown in Figure~\ref{fig:baselines}. Each applies a single encoder to a feature stream and classifies the result.

The multilayer perceptron is the simplest of these encoders, and it serves as the first baseline for both feature streams. An MLP accepts only a fixed-length vector, so the per-frame OpenFace descriptor is flattened into one long vector before the network classifies it (\texttt{openface\_mlp}), and the I3D clip vector, already fixed in length, is passed to an MLP directly (\texttt{i3d\_mlp}). Flattening discards the order of the frames, so this baseline reads the clip as an unordered collection of measurements and marks the level that any time-aware model has to improve upon.

The OpenFace descriptor is a sequence, however, and its order carries information that flattening throws away. A clip is an ordered series of frames in the same way that a sentence is an ordered series of words: engagement shows in how the face moves from one frame to the next, not in any single frame read on its own. Three encoders built to read ordered sequences are therefore applied to the raw $T\times709$ OpenFace stream --- a temporal convolutional network (\texttt{openface\_tcn}), a recurrent network (\texttt{openface\_lstm}), and a Transformer (\texttt{openface\_transformer}). Each stands for one of the main families of sequence models, and each preserves the frame order that the MLP removes.

The I3D stream cannot be read in the same way. I3D already averages a whole clip into a single vector when the feature is extracted, so no per-frame sequence remains for a temporal encoder to consume. Its single fixed-length vector fits only an MLP, which is why the I3D stream supports one baseline while the OpenFace stream supports four.

\begin{figure}[htbp]
\centering
\begin{subfigure}{\linewidth}\centering
\ifhanddrawn\includegraphics[width=\linewidth]{openface_mlp.png}\else\resizebox{\linewidth}{!}{%
\begin{tikzpicture}[node distance=6mm,>=Latex]
\node[abox](n1){Input\\$T\times709$};
\node[abox,right=of n1](n2){Flatten\\$T\cdot709$};
\node[abox,right=of n2](n3){Linear\\256};
\node[abox,right=of n3](n4){ReLU\\Dropout 0.3};
\node[abox,right=of n4](n5){Linear\\128};
\node[abox,right=of n5](n6){ReLU\\Dropout 0.3};
\node[abox,right=of n6](n7){Linear\\4};
\draw[->](n1)--(n2);\draw[->](n2)--(n3);\draw[->](n3)--(n4);\draw[->](n4)--(n5);\draw[->](n5)--(n6);\draw[->](n6)--(n7);
\end{tikzpicture}}\fi
\caption{\texttt{openface\_mlp}.}
\label{fig:openface_mlp}
\end{subfigure}\\[1.2ex]
\begin{subfigure}{\linewidth}\centering
\ifhanddrawn\includegraphics[width=\linewidth]{openface_tcn.png}\else\resizebox{\linewidth}{!}{%
\begin{tikzpicture}[node distance=6mm and 7mm,>=Latex]
\node[abox](n1){Input\\$T\times709$};
\node[abox,right=of n1](n2){TemporalBlock 1\\256 ch, dilation 1};
\node[abox,right=of n2](n3){TemporalBlock 2\\128 ch, dilation 2};
\node[abox,right=of n3](n4){TemporalBlock 3\\128 ch, dilation 4};
\node[abox,below=12mm of n1](m1){AdaptiveAvgPool\\128};
\node[abox,right=of m1](m2){Dropout 0.3};
\node[abox,right=of m2](m3){Linear\\128};
\node[abox,right=of m3](m4){ReLU\\Dropout 0.3};
\node[abox,right=of m4](m5){Linear\\4};
\draw[->](n1)--(n2);\draw[->](n2)--(n3);\draw[->](n3)--(n4);
\draw[->](n4.south) -- ++(0,-6mm) -| (m1.north);
\draw[->](m1)--(m2);\draw[->](m2)--(m3);\draw[->](m3)--(m4);\draw[->](m4)--(m5);
\end{tikzpicture}}\fi
\caption{\texttt{openface\_tcn}.}
\label{fig:openface_tcn}
\end{subfigure}\\[1.2ex]
\begin{subfigure}{\linewidth}\centering
\ifhanddrawn\includegraphics[width=\linewidth]{openface_lstm.png}\else\resizebox{\linewidth}{!}{%
\begin{tikzpicture}[node distance=6mm and 7mm,>=Latex]
\node[abox](n1){Input\\$T\times709$};
\node[abox,right=of n1](n2){LSTM layer 1\\hidden 256};
\node[abox,right=of n2](n3){Dropout 0.3};
\node[abox,right=of n3](n4){LSTM layer 2\\hidden 256};
\node[abox,below=12mm of n1](m1){Dropout 0.3};
\node[abox,right=of m1](m2){Linear\\128};
\node[abox,right=of m2](m3){ReLU\\Dropout 0.3};
\node[abox,right=of m3](m4){Linear\\4};
\draw[->](n1)--(n2);\draw[->](n2)--(n3);\draw[->](n3)--(n4);
\draw[->](n4.south) -- ++(0,-6mm) -| (m1.north);
\draw[->](m1)--(m2);\draw[->](m2)--(m3);\draw[->](m3)--(m4);
\end{tikzpicture}}\fi
\caption{\texttt{openface\_lstm}.}
\label{fig:openface_lstm}
\end{subfigure}
\caption{The five single-encoder baselines.}
\label{fig:baselines}
\end{figure}

\begin{figure}[htbp]\ContinuedFloat
\centering
\begin{subfigure}{\linewidth}\centering
\ifhanddrawn\includegraphics[width=\linewidth]{openface_transformer.png}\else\resizebox{\linewidth}{!}{%
\begin{tikzpicture}[node distance=6mm and 7mm,>=Latex]
\node[abox](n1){Input\\$T\times709$};
\node[abox,right=of n1](n2){Input projection\\Linear 128};
\node[aop,right=7mm of n2](add){$\oplus$};
\node[abox,above=5mm of add](pe){Positional encoding\\sinusoidal 128};
\node[abox,right=7mm of add](d1){EncoderLayer 1\\d 128, heads 4, FF 256};
\node[abox,right=of d1](d2){EncoderLayer 2\\d 128, heads 4, FF 256};
\node[abox,below=12mm of n1](m1){Mean-pool over $T$\\128};
\node[abox,right=of m1](m2){LayerNorm 128\\Dropout 0.3};
\node[abox,right=of m2](m3){Linear\\128};
\node[abox,right=of m3](m4){ReLU\\Dropout 0.3};
\node[abox,right=of m4](m5){Linear\\4};
\draw[->](n1)--(n2);\draw[->](n2)--(add);\draw[->](pe)--(add);\draw[->](add)--(d1);\draw[->](d1)--(d2);
\draw[->](d2.south) -- ++(0,-6mm) -| (m1.north);
\draw[->](m1)--(m2);\draw[->](m2)--(m3);\draw[->](m3)--(m4);\draw[->](m4)--(m5);
\end{tikzpicture}}\fi
\caption{\texttt{openface\_transformer}.}
\label{fig:openface_transformer}
\end{subfigure}\\[1.2ex]
\begin{subfigure}{\linewidth}\centering
\ifhanddrawn\includegraphics[width=\linewidth]{i3d_mlp.png}\else\resizebox{\linewidth}{!}{%
\begin{tikzpicture}[node distance=6mm,>=Latex]
\node[abox](n1){Input\\I3D clip 1024-d};
\node[abox,right=of n1](n2){Linear\\$1024\to256$};
\node[abox,right=of n2](n3){ReLU\\Dropout 0.3};
\node[abox,right=of n3](n4){Linear\\$256\to128$};
\node[abox,right=of n4](n5){ReLU\\Dropout 0.3};
\node[abox,right=of n5](n6){Linear\\$128\to4$};
\draw[->](n1)--(n2);\draw[->](n2)--(n3);\draw[->](n3)--(n4);\draw[->](n4)--(n5);\draw[->](n5)--(n6);
\end{tikzpicture}}\fi
\caption{\texttt{i3d\_mlp}.}
\label{fig:i3d_mlp}
\end{subfigure}
\caption{The five single-encoder baselines (continued).}
\end{figure}

\FloatBarrier
\section{Proposed Hybrid Model}
\label{sec:hybrid}

\subsection{Grouping the facial features by meaning}
\label{sec:groups}

The four OpenFace baselines share one weakness. A single encoder reads all 709 features at once, and so applies one inductive bias --- one built-in assumption about how a signal changes over time --- to every feature alike. No single assumption can fit them all, because the 709 features describe several distinct facial subsystems and those subsystems behave in very different ways: gaze and the eye landmarks shift in rapid, jittery micro-movements, head pose drifts slowly and smoothly, and the action units fire in short, sparse bursts. An encoder tuned to track quick eye movements is the wrong instrument for the slow drift of the head, so a single encoder over the whole descriptor suits a few of the features and misfits the rest.

Separating the features by the subsystem they describe removes this conflict. When a group holds only features that behave alike, it can be handed to the encoder whose assumption matches that behaviour, and no encoder is asked to serve incompatible dynamics at once. The 709 OpenFace features are accordingly split into five groups by meaning, matched from the column names in the fixed order used throughout the pipeline: gaze, eye landmarks, face landmarks, head pose, and action units. Table~\ref{tab:groups} lists the grouping, in which every one of the 709 columns belongs to exactly one group. This grouping by meaning is the central idea of the proposed model.

\begin{table}[ht]
\centering
\caption{Grouping of the 709-dimensional OpenFace descriptor into five feature groups~by~meaning.}
\label{tab:groups}
\begin{tabular}{llr}
\toprule
Group & Content & Dim \\
\midrule
Gaze          & gaze direction vectors and angles            & 8   \\
Eye landmarks & 2-D/3-D eye-region landmark coordinates        & 280 \\
Face landmarks& 2-D/3-D facial landmark coordinates            & 340 \\
Head pose     & head translation and rotation parameters      & 46  \\
Action units  & facial action-unit presence and intensity     & 35  \\
\midrule
\textbf{Total} & & \textbf{709} \\
\bottomrule
\end{tabular}
\end{table}

\FloatBarrier
\subsection{Proposed model architecture}

The proposed hybrid model is shown in Figure~\ref{fig:hybrid}. It keeps the three encoder families used by the OpenFace baselines --- the TCN, the LSTM, and the Transformer --- but instead of applying one of them to the whole $T\times709$ descriptor, it gives each of the five feature groups its own encoder, combines the results, and supervises every stream with its own auxiliary head. Reusing the same three families is a deliberate choice: it ensures that any gain over the baselines can be traced to the grouping alone, and not to a more powerful kind of encoder.

\begin{figure}[htbp]
\centering
\ifhanddrawn\includegraphics[width=\linewidth]{hybrid.png}\else\resizebox{\linewidth}{!}{%
\begin{tikzpicture}[>=Latex]
% inputs
\node[abox](gz){gaze\\$T\times8$};
\node[abox,below=3mm of gz](ey){eye\\$T\times280$};
\node[abox,below=3mm of ey](fc){face\\$T\times340$};
\node[abox,below=3mm of fc](hp){head pose\\$T\times46$};
\node[abox,below=3mm of hp](au){action units\\$T\times35$};
\node[abox,below=7mm of au](i3){I3D\\1024};
% encoders
\node[agrp,right=12mm of gz](gze){Group encoder\\64};
\node[agrp,right=12mm of ey](eye){Group encoder\\64};
\node[agrp,right=12mm of fc](fce){Group encoder\\64};
\node[agrp,right=12mm of hp](hpe){Group encoder\\64};
\node[agrp,right=12mm of au](aue){Group encoder\\64};
\node[abox,right=12mm of i3](i3e){Linear 128\\LayerNorm};
\draw[->](gz)--(gze);\draw[->](ey)--(eye);\draw[->](fc)--(fce);\draw[->](hp)--(hpe);\draw[->](au)--(aue);\draw[->](i3)--(i3e);
% concatenate (tall)
\node[abox,right=14mm of fce,minimum height=50mm,text width=16mm](cat){Concatenate\\$5\times64$\\$+\,128$\\$=448$};
\draw[->](gze.east)--(cat.west);\draw[->](eye.east)--(cat.west);\draw[->](fce.east)--(cat.west);
\draw[->](hpe.east)--(cat.west);\draw[->](aue.east)--(cat.west);\draw[->](i3e.east)-|(cat.south);
% fusion + output
\node[abox,right=10mm of cat](f1){Linear\\128};
\node[abox,right=6mm of f1](f2){ReLU\\Dropout 0.3};
\node[abox,right=6mm of f2](f3){Linear\\4};
\node[abox,right=6mm of f3](out){Ordinal level\\E0 / E1 / E2 / E3};
\draw[->](cat)--(f1);\draw[->](f1)--(f2);\draw[->](f2)--(f3);\draw[->](f3)--(out);
% auxiliary heads (one per stream)
\node[abox,above=9mm of cat](aux){Aux head per stream\\Linear 4};
\begin{scope}[on background layer]
\node[draw,dotted,rounded corners,fit=(gze)(eye)(fce)(hpe)(aue)(i3e),inner sep=2mm](encfit){};
\end{scope}
\draw[->,dashed](encfit.north)--(aux);
\end{tikzpicture}}\fi
\caption{Proposed hybrid model architecture (I3D stream enabled; the I3D-stream-disabled variant drops the I3D row).}
\label{fig:hybrid}
\end{figure}

Each group is read by its own \texttt{GroupEncoder}, one of the three interchangeable encoders drawn in Figure~\ref{fig:group_encoders}. Whichever encoder a group is assigned, its per-frame outputs are averaged over time and layer-normalised, so that every group leaves a 64-dimensional embedding of the same width regardless of how many raw features it began with. Fixing one common width matters because the groups differ greatly in size, from eight features for gaze to 340 for the face landmarks, and a shared width lets each contribute to the fused representation on an equal footing rather than in proportion to its raw dimensionality.

\begin{figure}[htbp]
\centering
\begin{subfigure}{\linewidth}
\centering
\ifhanddrawn\includegraphics[width=\linewidth]{hybrid_tcn_encoder.png}\else\resizebox{\linewidth}{!}{%
\begin{tikzpicture}[node distance=6mm,>=Latex]
\node[abox](n1){Input\\$T\times D_{\text{group}}$};
\node[abox,right=of n1](n2){TemporalBlock 1\\64 ch, dilation 1};
\node[abox,right=of n2](n3){TemporalBlock 2\\64 ch, dilation 2};
\node[abox,right=of n3](n4){Mean-pool\\64};
\node[abox,right=of n4](n5){LayerNorm\\64};
\draw[->](n1)--(n2);\draw[->](n2)--(n3);\draw[->](n3)--(n4);\draw[->](n4)--(n5);
\end{tikzpicture}}\fi
\caption{TCN group encoder.}
\label{fig:ge_tcn}
\end{subfigure}\\[2ex]
\begin{subfigure}{\linewidth}
\centering
\ifhanddrawn\includegraphics[width=\linewidth]{hybrid_transformer_encoder.png}\else\resizebox{\linewidth}{!}{%
\begin{tikzpicture}[node distance=6mm and 7mm,>=Latex]
\node[abox](n1){Input\\$T\times D_{\text{group}}$};
\node[abox,right=of n1](n2){Input projection\\Linear $D_{\text{group}}\to64$};
\node[aop,right=7mm of n2](add){$\oplus$};
\node[abox,above=5mm of add](pe){Positional encoding\\sinusoidal 64-d};
\node[abox,right=7mm of add](d1){EncoderLayer 1\\d 64, heads 4, FF 128};
\node[abox,right=of d1](d2){EncoderLayer 2\\d 64, heads 4, FF 128};
\node[abox,right=of d2](p){Mean-pool over $T$\\64};
\node[abox,right=of p](ln){LayerNorm\\64};
\draw[->](n1)--(n2);\draw[->](n2)--(add);\draw[->](pe)--(add);\draw[->](add)--(d1);\draw[->](d1)--(d2);\draw[->](d2)--(p);\draw[->](p)--(ln);
\end{tikzpicture}}\fi
\caption{Transformer group encoder.}
\label{fig:ge_transformer}
\end{subfigure}\\[2ex]
\begin{subfigure}{\linewidth}
\centering
\ifhanddrawn\includegraphics[width=\linewidth]{hybrid_lstm_encoder.png}\else\resizebox{\linewidth}{!}{%
\begin{tikzpicture}[node distance=6mm,>=Latex]
\node[abox](n1){Input\\$T\times D_{\text{group}}$};
\node[abox,right=of n1](n2){LSTM layer 1\\hidden 64};
\node[abox,right=of n2](n3){Dropout 0.2};
\node[abox,right=of n3](n4){LSTM layer 2\\hidden 64};
\node[abox,right=of n4](n5){Mean-pool over $T$\\64};
\node[abox,right=of n5](n6){LayerNorm\\64};
\draw[->](n1)--(n2);\draw[->](n2)--(n3);\draw[->](n3)--(n4);\draw[->](n4)--(n5);\draw[->](n5)--(n6);
\end{tikzpicture}}\fi
\caption{LSTM group encoder.}
\label{fig:ge_lstm}
\end{subfigure}
\caption{The three interchangeable per-group encoders.}
\label{fig:group_encoders}
\end{figure}

The five 64-dimensional embeddings are concatenated into a single 320-dimensional vector, which a two-layer classifier ($320\!\to\!128\!\to\!4$) maps to the four engagement logits. Each group also feeds a small auxiliary classifier of its own, and the training objective adds a fraction of the averaged auxiliary losses --- weighted at $0.2$ --- to the main loss (Section~\ref{sec:losses-method}). These auxiliary heads keep every stream discriminative on its own, so that no single group can be ignored while the others carry the prediction.

An optional sixth stream carries the I3D feature. Because I3D supplies a single pooled vector for the whole clip rather than a per-frame sequence, it is read by an MLP ($1024\!\to\!256\!\to\!128$) instead of a temporal encoder, and layer-normalised to a 128-dimensional embedding with its own auxiliary head. Enabling this stream widens the fused vector to 448 dimensions ($5\times64+128$); disabling it leaves the model on the five OpenFace groups alone, and both variants are carried forward for evaluation.

A single configuration of the model is therefore fixed by the encoder chosen for each of the five groups, written as an \texttt{arch\_key} of five tokens in group order --- for example \texttt{TCN\_T\_TCN\_LSTM\_T} (\texttt{T}~=~Transformer, \texttt{TCN}~=~TCN, \texttt{LSTM}~=~LSTM). With three encoders available to each group this gives $3^{5}=243$ possible configurations, and each can run with the I3D stream enabled or disabled.

\end{document}
